\begin{tuluyen}
Trong một lon cà phê sử dụng công nghệ tự làm nóng, ngăn chứa dung dịch cà phê được tách riêng với ngăn chứa CaO và ngăn chứa nước. Khi sử dụng, người dùng kích hoạt để trộn lẫn CaO và nước. Nhiệt tỏa ra của phản ứng sẽ làm nóng ngăn chứa dung dịch cà phê.

\begin{enumerate}
\item Viết phương trình hóa học của CaO và nước.
\item Tính lượng CaO tối thiểu để làm nóng 124 mL dung dịch cà phê thêm 40$^\circ$C. Biết: 1 mol CaO phản ứng hết với nước tỏa ra nhiệt lượng là 62 kJ, khối lượng riêng của dung dịch cà phê là $d=1,0$ g/mL, nhiệt dung riêng của dung dịch cà phê là $C=4,18$ J/($^\circ$C.g). Nhiệt lượng cần cung cấp được tính theo công thức: $Q=m.C.\Delta T$, với $m$ là khối lượng dung dịch, $C$ là nhiệt dung riêng của dung dịch và $\Delta T$ là khoảng nhiệt độ cần tăng.
\end{enumerate}
\source{Trích nguyên văn Câu I.1, đề tuyển sinh vào lớp 10 Trường THPT Chuyên KHTN Hà Nội 2026.}
\dapso{(a) \ce{CaO + H2O -> Ca(OH)2}; (b) \(\approx\SI{18.73}{\gram}\) CaO.}
\begin{loigiaiBT}
Phương trình hóa học:
\[
\ce{CaO + H2O -> Ca(OH)2}.
\]
Khối lượng dung dịch cà phê:
\[
m=124\times1.0=\SI{124}{\gram}.
\]
Nhiệt lượng cần cung cấp:
\[
Q=124\times4.18\times40=\SI{20732.8}{\joule}=\SI{20.7328}{\kilo\joule}.
\]
Số mol CaO cần phản ứng:
\[
n_{\ce{CaO}}=\frac{20.7328}{62}\approx\SI{0.3344}{\mole}.
\]
Khối lượng CaO:
\[
m_{\ce{CaO}}=0.3344\times56\approx\SI{18.73}{\gram}.
\]
\textbf{Vậy,} lượng CaO tối thiểu cần dùng khoảng \SI{18.73}{\gram}.
\end{loigiaiBT}
\end{tuluyen}
